% Generated draft: requires mathematical and source review.
\documentclass[11pt]{article}
\usepackage[margin=25mm]{geometry}
\usepackage{fontspec}
\setmainfont{Latin Modern Roman}
\usepackage{amsmath,amssymb,mathtools,unicode-math}
\setmathfont{Latin Modern Math}
\usepackage{enumitem,needspace,xurl,hyperref}
\setlength{\parindent}{0pt}
\setlength{\parskip}{5pt}
\setlength{\emergencystretch}{3em}
\newcommand{\N}{\mathbb{N}}
\newcommand{\Z}{\mathbb{Z}}
\newcommand{\Q}{\mathbb{Q}}
\newcommand{\R}{\mathbb{R}}
\newcommand{\C}{\mathbb{C}}
\newcommand{\F}{\mathbb{F}}
\newcommand{\A}{\mathbb{A}}
\newcommand{\PP}{\mathbb{P}}
\newcommand{\E}{\mathbb{E}}
\newcommand{\eps}{\varepsilon}
\newcommand{\dd}{\,\mathrm{d}}
\begin{document}
\section*{UnsolvedMath ID: 8800096}
\textbf{Source problem code:} catalogue entry
\subsection{Definitions and mathematical statement}
The supplied task is to identify and solve tractable special cases of a ``signature problem'' described in the original slides. That problem is not reproduced or otherwise identified in the record.

No input objects, meaning of signature, admissible outputs, equations, equivalence relation, decision predicate, or background hypotheses are supplied. Therefore a mathematical relation or algorithmic problem for the requested special cases cannot be reconstructed. In particular it is not justified to assume this is a digital-signature construction merely from the workshop's cryptography context.

A precise special-case task first needs the original problem and its input encoding, then a specified restriction and tractability criterion, such as a complexity or resource guarantee. The source does not choose any of these. This draft retains the request for both identifying special classes and solving them, while recording that no standalone quantified formulation is possible from the supplied evidence.
\subsection{Short English statement}
Find solvable, computationally tractable cases of the signature problem in Huang's slides. The actual problem and the meaning of signature are missing. The supplied summary likewise says the specific problem could not be identified with certainty and reports no verified tractable-case resolution. Its OPEN-TRIAGE label is retained without a new status claim.
\subsection{Sources}
\begin{itemize}
\item Local UnsolvedMath registry, permanent ID 8800096, source code catalogue entry, ``Tractable special cases of the signature problem,'' proposed by Ming-Deh Huang. Dataset curation and metadata: Ulam AI and UnsolvedMath Contributors, CC BY 4.0.
\item IPAM, \emph{Workshop on Number Theory and Cryptography --- Open Problems} (2006), Attachment 6673, PDF page 30, schedule locator \url{https://www.ipam.ucla.edu/programs/workshops/workshop-i-number-theory-and-cryptography-open-problems/?tab=schedule}; pdf-and-legacy-ppt visual-audit extraction, accessed 2026-07-29.
\item The antecedent signature problem and any exact primary resolution reference are absent from the supplied record.
\item Source provenance reports inspection of 17 PDFs and two legacy PPT decks among 19 surviving attachments, excluding repetitions and solved definitions; this is supplied provenance, not a new inspection claimed here.
\item Current primary-literature status audit and redistribution terms NEEDS\_REVIEW for the publicly accessible author/publisher copy, separately from the dataset licence.
\end{itemize}
\end{document}
